\section{Vinculación con la Problemática Inicial}

En el Capítulo \ref{chap:introduccion} se establecieron tres problemas de partida que justifican el desarrollo de la plataforma: Primero, la dificultad de reacción del inversor minorista ante un mercado cripto continuo y altamente volátil. Segundo, la barrera técnica para adoptar trading algorítmico de forma segura. Tercero, el sesgo emocional en la toma de decisiones de inversión.

Este capítulo se estructura para responder de forma explícita a dichos problemas. En primer lugar, la arquitectura por capas y el orquestador Java abordan el segundo problema al desacoplar responsabilidades y reducir la complejidad operativa del sistema para su uso académico. En segundo lugar, el pipeline ETL, la persistencia masiva y la gestión de concurrencia permiten tratar el primer problema, al sostener la ingesta continua de datos y el procesamiento oportuno de velas de mercado. Por último, los módulos de backtesting, \textit{paper trading} e IA se orientan a mitigar el tercer problema, al desplazar decisiones discrecionales hacia reglas reproducibles y validables con evidencia empírica.
